\section{The Anatomy of the Dynamic Gap}
\label{sec:anatomy}

Section~\ref{sec:separation} establishes that dynamic count-valued testing can
strictly outperform static designs. This section asks a finer question: what is
that advantage made of, and how much of it survives the structural restrictions
that make the problem computationally tractable?

We decompose the gap along two orthogonal axes --- the \emph{observation
channel} (binary versus count-valued) and the \emph{design mode} (static versus
adaptive) --- and then measure separately the price of restricting attention to
laminar pool histories. All quantities reported here are exact: they are
obtained by enumerating the belief state, not by sampling.

\paragraph{Conventions.}
Throughout, $X_i=1$ indicates that agent $i$ is infected, $p_i=\Pr[X_i=1]$, and
$q_i=1-p_i$ is the prior probability that agent $i$ is healthy. The regime of
interest is $q_i<\tfrac12$, in which healthy agents are the scarce resource. We
write $r(t,X)=\sum_{i\in t}X_i$ for the count returned by an augmented test on
pool $t$.

Two coarsenings of the count are natural, and they are not equivalent:
\[
  Y^{\mathrm{inf}}(t)=\mathbf{1}\{r(t,X)>0\},
  \qquad
  Y^{\mathrm{heal}}(t)=\mathbf{1}\{r(t,X)<|t|\}.
\]
The first is the classical group-testing channel and reports the presence of an
infected agent; the second reports the presence of a \emph{healthy} agent. In a
welfare-maximizing setting the object of interest is the healthy population, and
$Y^{\mathrm{heal}}$ is the binary channel we study. Every binary quantity below
refers to $Y^{\mathrm{heal}}$.

We also distinguish two certification rules, and the distinction turns out to be
consequential rather than cosmetic.

\begin{definition}[Certification rules]
\label{def:certification}
Fix a testing history $H$ with observed outcomes. Under \emph{credited-only}
certification, agent $i$ contributes $u_i$ to welfare if and only if $i$ belongs
to some tested pool $t$ with observed count $r(t)=0$. Under
\emph{inference-based} certification, agent $i$ contributes $u_i$ if and only if
$X_i=0$ in every assignment consistent with $H$.
\end{definition}

Credited-only certification is the hard-clearing convention: a deduction informs
future decisions but does not itself generate reward. Inference-based
certification counts an agent whenever the observed constraint system determines
its status. Every agent credited under the first rule is certified under the
second, but not conversely.

\subsection{Two Orthogonal Ingredients}
\label{sec:anatomy-ladder}

We first quantify the two ingredients separately. Consider $N=10$ homogeneous
agents with $q_i=0.2$ and $u_i=1$, no pool-size constraint, inference-based
certification, and a testing budget $B$. Table~\ref{tab:ladder} reports the exact
optimum of each class, obtained by complete enumeration of the belief state.

\begin{table}[h]
\centering
\begin{tabular}{cccccc}
\toprule
$B$ & static binary & dynamic binary & static count & laminar count & dynamic count \\
\midrule
1 & 0.200 & 0.200 & 0.200 & 0.200 & 0.200 \\
2 & 0.400 & 0.488 & 0.408 & 0.536 & 0.536 \\
3 & 0.600 & 0.790 & 0.800 & 0.928 & 1.000 \\
\bottomrule
\end{tabular}
\caption{Exact expected welfare by class, $N=10$, $q_i=0.2$, $u_i=1$,
inference-based certification. ``Laminar count'' restricts the dynamic
count-valued optimum to histories in which any two pools are nested or
disjoint.}
\label{tab:ladder}
\end{table}

Three observations follow. First, at $B=1$ all five classes coincide, which fixes
the base case. Second, at $B=3$ the two ingredients are of comparable value in
isolation: counting without adaptivity yields $0.800$ and adaptivity without
counting yields $0.790$, an increase of $33\%$ and $32\%$ respectively over the
static binary baseline. Neither dominates the other. Third, the two together
yield $1.000$, an increase of $67\%$, so the ingredients are complementary rather
than redundant.

That counting alone should be worth as much as adaptivity alone is not obvious,
and the next subsection explains where the static count-valued gain comes from.

\subsection{Certification Rules and the Static Baseline}
\label{sec:anatomy-certification}

The value of the static count-valued class depends entirely on which
certification rule is in force. Under credited-only certification, singleton
designs are optimal in the high-prevalence regime.

\begin{proposition}[Optimality of singleton designs under credited-only certification]
\label{prop:singleton-optimal}
Let all agents be homogeneous with $q_i=q\le\tfrac12$ and $u_i=u$, and adopt
credited-only certification. Then no static design of $B$ tests achieves
expected welfare greater than $B\,u\,q$, and $B$ distinct singleton tests attain
this value.
\end{proposition}

\begin{proof}
Let $t_1,\ldots,t_B$ be any static design. An agent is credited only if it lies
in a pool whose observed count is zero, so by a union bound over pools the
expected number of credited agents is at most
$\sum_{j=1}^{B}|t_j|\Pr[r(t_j,X)=0]=\sum_{j=1}^{B}|t_j|\,q^{|t_j|}$.
For $g\ge1$ the ratio of consecutive terms of $g\mapsto g\,q^{g}$ is
$q(g+1)/g\le 2q\le 1$, so the sequence is nonincreasing and
$\max_{g\ge1} g\,q^{g}=q$. Hence the expected welfare is at most $B\,u\,q$.
Singleton tests on distinct agents attain it, since each is credited
independently with probability $q$.
\end{proof}

Under inference-based certification the conclusion fails, and it fails for a
structural reason: a count is an integer-valued linear measurement, and a small
number of such measurements can identify more agents than they test
individually.

\begin{proposition}[Three count-valued tests identify four agents]
\label{prop:coin-weighing}
There is a static design of three count-valued tests on four agents that
determines $X\in\{0,1\}^{4}$ exactly, for every realization. Consequently, under
inference-based certification with homogeneous parameters, the static
count-valued optimum with $B=3$ is at least $4uq$, which strictly exceeds the
bound $3uq$ of Proposition~\ref{prop:singleton-optimal}.
\end{proposition}

\begin{proof}
Label the agents $a,b,c,d$ and take
\[
  t_1=\{a,b,d\},\qquad t_2=\{a,c,d\},\qquad t_3=\{b,c,d\},
\]
so that each agent belongs to a distinct subset of the tests. Writing
$r_1,r_2,r_3$ for the observed counts,
\[
  r_1+r_2+r_3 \;=\; 2(X_a+X_b+X_c)+3X_d,
\]
whose parity equals the parity of $X_d$; since $X_d\in\{0,1\}$, this determines
$X_d$. Given $X_d$, the three equations become
$X_a+X_b=r_1-X_d$, $X_a+X_c=r_2-X_d$, $X_b+X_c=r_3-X_d$, whose sum is
$2(X_a+X_b+X_c)$ and therefore determines $X_a+X_b+X_c$; subtracting each
pairwise sum recovers each variable individually. All four agents are thus
determined, and the $4q$ of them that are healthy in expectation are certified.
\end{proof}

The same statement is impossible for binary tests, by a counting argument that
requires no design at all.

\begin{proposition}[Binary channels cannot identify four agents with three tests]
\label{prop:binary-impossible}
No static design of three binary tests determines $X\in\{0,1\}^{4}$ for every
realization.
\end{proposition}

\begin{proof}
Three binary tests induce at most $2^{3}=8$ distinct outcome vectors, while
there are $2^{4}=16$ realizations. By pigeonhole two distinct realizations
produce the same outcome vector; they differ in the status of some agent, which
therefore cannot be determined.
\end{proof}

Proposition~\ref{prop:coin-weighing} is an instance of the coin-weighing problem
of Erd\H{o}s and R\'enyi: $B$ measurements that report subset sums identify on
the order of $B\log B$ binary unknowns. Under inference-based certification the
static count-valued optimum therefore grows superlinearly in the budget, rather
than as $B\,u\,q$. Any separation argument that measures the dynamic advantage
against $B\,u\,q$ is correct under credited-only certification and compares
against the wrong baseline under the inference-based rule. We adopt
credited-only certification as the normative convention of the paper and record
the alternative explicitly, since the choice determines the value of an entire
class.

\subsection{The Cost of Laminarity}
\label{sec:anatomy-laminar}

Laminar histories admit exact inference (Section~\ref{sec:laminar-pools}), and
with count-valued tests the structure is even stronger: every tested pool ends
with a known count, so the tested population decomposes into residual atoms with
known counts and the untested territory retains its product prior.

\begin{proposition}[Whole atoms carry no information]
\label{prop:atoms-constant}
Let $H$ be a laminar history of count-valued tests with residual atoms
$\{D_t\}$ and residual counts $\{c(D_t)\}$, and let $V(H)$ denote the untested
territory. Let $s$ be a new pool that is the union of a subfamily
$\mathcal{A}$ of whole residual atoms with a subset $v\subseteq V(H)$. Then
\[
  r(s,X) \;=\; \sum_{D\in\mathcal{A}} c(D) \;+\; r(v,X),
\]
so the observation $r(s,X)$ is informationally equivalent to $r(v,X)$ alone.
\end{proposition}

\begin{proof}
By Lemma~\ref{lem:count-equivalence}(a) the history determines $r(D,X)=c(D)$ for
every residual atom $D$. The atoms are pairwise disjoint and disjoint from
$V(H)$ by Lemma~\ref{lem:atomic-partition}, so counts add over the blocks of
$s$. The first summand is a constant given $H$, and subtracting a constant from
an observation preserves the induced partition of the sample space.
\end{proof}

Proposition~\ref{prop:atoms-constant} has an algorithmic consequence: when
searching over laminar continuations it suffices to consider fresh pools and
strict subsets of a single atom, since enlarging a pool by whole atoms changes
nothing. This reduces the laminar state to a multiset of (size, count) pairs
together with the size of the untested territory, and it is the reason the
laminar optimum remains computable at budgets where the unrestricted optimum
does not: in our implementation the laminar solver reaches $B=8$ in seconds,
while the unrestricted solver does not terminate at $B=4$.

Tractability, however, is not free. Table~\ref{tab:ladder} shows that at $B=3$
the laminar optimum is $0.928$ against an unrestricted optimum of $1.000$, a
relative loss of $7.2\%$, while at $B\le2$ the restriction is costless. The
optimal unrestricted policy crosses a previously tested pool, and the mechanism
is worth stating explicitly.

\begin{example}
\label{ex:crossing}
Consider $N=10$ homogeneous agents with $q_i=0.2$, $u_i=1$, and $B=3$ under
inference-based certification. The optimal policy first tests a pool
$\{a,b,c\}$. On the outcome $r=2$, exactly one of the three agents is healthy
and each has posterior health probability $\tfrac13$. The optimal second test is
$\{a,b,d\}$, where $d$ is untested --- a pool that neither contains, is contained
in, nor is disjoint from $\{a,b,c\}$.

The crossing test performs two functions simultaneously. If it returns a count
of three, then $a$, $b$ and $d$ are all infected, so $c$ is healthy and is
certified without ever being tested, and the remaining budget is not spent on
$d$. If it returns a count of one, then exactly one of $\{a,b\}$ is healthy and
so is $d$, which is certified; the last test resolves the pair. A laminar policy
must choose between probing the tested group and probing fresh territory; the
crossing test uses the residual information of the group to extract value from a
new agent with the same test.
\end{example}

The loss reported here is a single instance and is evidence, not a bound;
adversarial search over the same family has located instances with a ratio of
$0.9069$.

\subsection{Where Policies Leave the Laminar Class}
\label{sec:anatomy-falsifier}

Example~\ref{ex:crossing} raises the question of how often crossing matters. We
answer it by classifying every decision taken by a policy at every reachable
state, weighted by the probability of reaching that state.

\begin{definition}[Crossing]
\label{def:crossing}
A pool $t$ \emph{crosses} a pool $s$ if $t\cap s\notin\{\varnothing,t,s\}$. An
action is crossing with respect to a history $H$ if it crosses some pool already
executed in $H$. Disjoint pools, descendants, ancestors and repetitions are not
crossing.
\end{definition}

For a policy $\pi$ and a class $\kappa$ of actions we write
$W^{\pi}_{\kappa}=\sum_{H}\Pr^{\pi}[H]\,\mathbf{1}\{\pi(H)\in\kappa\}$ for the
probability mass of decisions of that class. Table~\ref{tab:crossing} reports
$W^{\pi}_{\mathrm{cross}}$ over a sweep of $54$ instances with
$N\in\{4,5,6\}$, $B\in\{1,2,3\}$, pool-size constraint $G\in\{2,3\}$ and
$q\in\{0.05,0.45,0.90\}$, comprising $872$ decisions in total.

\begin{table}[h]
\centering
\begin{tabular}{lcc}
\toprule
Policy & crossing mass & mean local regret \\
\midrule
myopic $S_0$ & 0.0026 & 0.0160 \\
rollout & 0.0215 & 0.0000 \\
unrestricted optimum & 0.0367 & 0.0026 \\
\bottomrule
\end{tabular}
\caption{Probability mass of crossing decisions, and mean local regret against
the rollout $Q$-function, averaged over $54$ instances.}
\label{tab:crossing}
\end{table}

The crossing mass is strictly increasing in the quality of the policy: the
optimum crosses roughly fourteen times as often as the myopic policy. The
laminar restriction therefore binds precisely on the policies that are worth
restricting, which is the empirical counterpart of Example~\ref{ex:crossing}.

The second column records a further observation. The rollout policy maximizes
$Q^{\mathrm{g}}$ by construction and so has zero local regret; the unrestricted
optimum does not, because $Q^{\mathrm{g}}$ evaluates an action under a greedy
continuation rather than an optimal one. The magnitude of that regret quantifies
the residual between rollout and optimality.

\subsection{Planning without Solving}
\label{sec:anatomy-rollout}

Exact dynamic optimization is infeasible beyond small instances, which motivates
policies that plan without solving. The rollout policy scores each action by its
value under a greedy continuation,
\[
  Q^{\mathrm{g}}_b(H,t)
  =\mathbb{E}\!\left[r(H,t,R)+V^{\mathrm{g}}_{b-1}(H')\,\middle|\,H,t\right],
\]
executes the maximizer, and re-plans at the next state rather than continuing
greedily. Proposition~\ref{prop:policy-improvement} guarantees
$V^{\mathrm{r}}\ge V^{\mathrm{g}}$ at every reachable state.

We verify the implementation rather than assume it. The hypothesis that the
greedy action belongs to the rollout candidate set is checked state by state,
as is the dominance conclusion. Expected values are computed by two independent
procedures: a backward recursion over the belief tree, which branches on
observed counts and weights by conditional probabilities, and a forward
enumeration over all $2^{N}$ latent profiles, which runs the policy
deterministically on each profile and averages under the prior. No probability
appears inside the second computation. Over the test battery the two agree to
within $2.22\times10^{-16}$.

On the instance $N=5$, $q_i=0.3$, $u_i=1$, $B=3$, $G=2$, the myopic policy
attains $0.900$ and the rollout attains $1.011$, an improvement of $12.3\%$. The
myopic value coincides digit for digit with the static optimum, since a myopic
policy on homogeneous parameters never conditions on what it observes and is
therefore a fixed design in disguise.

\paragraph{A diagnostic for surrogate objectives.}
The same machinery supplies a falsification test for candidate scoring
functions. In the family with homogeneous parameters and small $q$, the plan
that covers the population with $k$ disjoint pools of size $G$ and then locates a
healthy agent by binary search attains, under credited-only certification,
\[
  u\left(1-(1-q)^{kG}\right),
  \qquad k=\max\{0,\;B-\lceil\log_2 G\rceil-1\},
\]
where the subtracted test is the one that converts a located agent into credited
utility. At $(q,G,k,B)=(0.05,16,2,7)$ this gives $0.8063u$ against a singleton
baseline of $0.35u$, a ratio of $2.30$. A scoring function that does not value
this plan cannot recover the separation, whatever its form. The myopic score
$S_0$ fails exactly two of the nine trajectory checks: it does not value opening
fresh territory, since a pool of size $g$ is scored by $q^{g}$, and it is blind
to the remaining budget. These are precisely the two properties a planning
score must supply.

\subsection{Open Questions}
\label{sec:anatomy-open}

\begin{enumerate}
  \item Is the near-equality between the value of counting alone and the value
    of adaptivity alone a coincidence of these parameters, or does it reflect a
    structural relation between the two relaxations?

  \item Under inference-based certification, what is the exact growth rate of the
    static count-valued optimum, and does the coin-weighing construction remain
    optimal when utilities are heterogeneous?

  \item Proposition~\ref{prop:singleton-optimal} is stated for homogeneous
    parameters. Does the singleton bound survive heterogeneous $q_i$ and $u_i$,
    where a pool contributes $\left(\prod_{i\in t}q_i\right)\sum_{i\in t}u_i$?

  \item Can the laminar class be relaxed to admit a bounded number of crossings
    while retaining exact inference, and does the resulting class recover the
    $7.2\%$?

  \item Given that the unrestricted optimum does not maximize $Q^{\mathrm{g}}$,
    is there a continuation policy, cheaper than optimal, whose induced
    $Q$-function is maximized by the optimum?

  \item Does a budget-aware surrogate exist that passes all nine trajectory
    checks while remaining strictly cheaper to evaluate than the rollout oracle?
\end{enumerate}

Taken together, these results suggest that the advantage of dynamic count-valued
testing is not a single phenomenon but the composition of three separable ones:
the resolution supplied by integer-valued observations, the conditioning
supplied by adaptivity, and the freedom to test pools that cut across the
history. The first two are of comparable magnitude and combine
superadditively; the third is small in mass but concentrated on exactly the
decisions that distinguish good policies from mediocre ones.
